\textbf{Motion of a charged ball}

A solid, homogeneous spherical ball of mass $m$ and radius $R$ is made of insulating material and has charge $Q$ distributed uniformly throughout its volume. The ball is placed on a large horizontal surface, and set in rolling motion without slipping in such a way that its center starts to move with initial horizontal velocity $v_0$. There is a uniform magnetic field (flux density) of magnitude $B$ perpendicular to the surface. The coefficient of static friction is large enough to prevent the ball from slipping on the surface. The moment of inertia of the ball about an axis through its center is $\frac{2mR^2}{5}$.

Describe the motion of the center of the ball and the shape of its trajectory.

Hint: Depending on your approach you may use the following identity:
\[
\vec{a}\ \mathrm{x}\left(\vec{b}\ \mathrm{x}\ \vec{c}\right)=\vec{b}\left(\vec{a}\cdot\vec{c}\right)-\vec{c}\left(\vec{a}\cdot\vec{b}\right)
\]
valid for any three vectors $\vec{a},\vec{b}$ a$nd\ \vec{c}$.
% NOTE: "and" is printed with an upright "a" and italic "nd" (typesetting accident in the source)
